\section{检验设计依据与适配效果}
\label{sec:evidence}

以下在 CLIP ViT-B/32 八任务的归一化 KL-G 路径上，固定专家、图像与修正规则，以 seed 0、每轮 CG100、最多四轮进行配对诊断。实验依次检验状态依赖、重估价值和起点适配；完整协议与对照见附录~\ref{app:protocol}。

\paragraph{统计是否随网络状态改变？}
Average 首轮前后，实际求解度量的余弦中位数为 $0.553$，固定诊断图像上的纠正梯度余弦为 $0.248$。图~\ref{fig:motivation}a 显示前期度量变化明显，后期趋稳。跨 block 干预进一步固定目标块权重为初始值，仅将周围网络替换为第一轮状态：Average 的 block 11 纠正梯度余弦仍降至 $0.4915$。这表明统计变化可以来自周围网络，不能只根据本块参数是否改变来判断是否需要适配。

\paragraph{处理这种变化是否有实际价值？}
固定同一个 Average 第一轮 checkpoint，使用当前与初始统计构造并接受更新，独立诊断 KL 分别降低 $0.080560$ 与 $0.044534$。进一步共享第一轮模型后分叉，四轮刷新与冻结统计的测试准确率为 $85.312\%$、$84.519\%$；共同第一轮为 $83.712\%$，单轮 CG400 为 $84.565\%$（图~\ref{fig:motivation}b）。这些对照区分继续更新与重估的额外收益。刷新收益依赖状态，RegMean 上仅为 $0.003$ 个百分点；CG400 也不是已收敛或等耗时对照。等幅原始提案和跨 block 的负结果见附录，现有证据不意味着当前统计在每个方向、每一轮都更好。

\paragraph{能否适配不同合并起点？}
Average/RegMean 的诊断准确率差距由 $22.36$ 缩至 $0.15$ 个百分点（图~\ref{fig:motivation}c），预测 JS 由 $0.13847$ 降至 $0.00412$；终态完整测试仅差 $0.24$ 个百分点，但逐层相对参数距离中位数仍为 $2.53\%$。这支持功能上的适配，而不要求参数同解。适配也有边界：Iso-C 与增量范数匹配的未拉平起点在四轮后仍相差 $8.25$ 个测试准确率百分点（表~\ref{tab:starts}）。

\begin{figure}[t]
 \centering
 \includegraphics[width=\linewidth]{motivation.pdf}
 \caption{\textbf{状态依赖、重估收益与起点适配。}ViT-B/32，八任务，seed 0。
 (a) 相邻学生状态的实际求解度量 $\bar A\otimes\bar G$ 的余弦距离，取层/任务中位数，对数纵轴。
 (b) Average 的完整测试准确率；刷新与冻结分支共享第一轮模型，均继续求解和选步，虚线为单轮 CG400。
 (c) 三种常用起点在同一独立诊断集上的准确率。(a) 比较校准统计，(b)/(c) 使用不同数据分割，数值不跨面板相减。}
 \label{fig:motivation}
\end{figure}
